\section{Conclusion}
\label{sec:conclusion}

This paper introduced the SAGE framework, a hierarchical system for evaluating literary quality through systematic decomposition into interpretive dimensions assessed by large language models. Through comprehensive evaluation of 100 short stories across three analytical layers using dual-mode assessment with multi-round iterative reflection, we address four research questions and demonstrate that LLM-based literary evaluation can achieve exceptional reliability while revealing systematic patterns distinguishing canonical literature, commercial genre fiction, and AI-generated narratives.

Regarding reliability (RQ1), the framework achieves robust performance with 100\% evaluation success, 98.8\% score convergence, and greater than 94\% inter-rater agreement, validating the multi-round architecture as an effective mechanism for grounding interpretive judgments. Regarding discriminative capacity (RQ2), the consistent genre hierarchy (Canonical $>$ Pulp $>$ LLM, all $p<0.001$) with layer-specific discrimination patterns (very large effect sizes for cultural and existential dimensions at $d>2.4$, smaller but still large effects for emotional dimension at $d$=1.68) reveals a capability profile distinguishing pattern-reproducible from stance-requiring literary capacities. Regarding evaluation robustness (RQ3), near-perfect mode invariance (maximum $|\Delta|$=0.05) demonstrates that content-based textual analysis and metadata-based critical synthesis converge on equivalent assessments. Regarding dimensional independence (RQ4), moderate cross-layer correlations ($r$ = 0.649--0.683) confirm that cultural, emotional-psychological, and existential-philosophical representation capture empirically distinguishable quality dimensions providing complementary rather than redundant information.

This work contributes to computational literary studies and LLM evaluation research in three respects. First, we provide a theoretically grounded framework operationalizing abstract critical concepts, from Bourdieu's cultural capital to Heideggerian existential analysis, through structured prompting strategies, validated by successful discrimination across quality categories. Second, we demonstrate that LLM-based evaluation with multi-round reflection, independent validation, and explicit theoretical grounding can achieve reliability comparable to human literary criticism, suggesting promise for LLM-assisted critical analysis beyond surface feature extraction. Third, we characterize systematic limitations of current generative models through controlled quality comparison, revealing that approximately 35\% performance gaps in cultural critique and philosophical depth indicate fundamental challenges in AI literary generation that surface fluency and emotional mimicry cannot bridge.

Several limitations warrant acknowledgment. The evaluation corpus comprises 100 English-language short stories, limiting generalizability to other languages, literary forms, and cultural contexts. The study employs a single LLM model (GPT-5-mini); cross-model comparison would distinguish framework-level reliability from model-specific behaviors. The reported inter-rater agreement (greater than 94\%) reflects consistency between two LLM-based evaluators (iterative and validator) rather than agreement with human expert judgments; validation against human literary critics remains an important direction for establishing external validity. The framework validates only interpretive layers (L4--L6), leaving integration with computational metrics (L1--L3) for future work. Promising extensions include expanding evaluation to non-English literatures, systematic cross-model comparison across LLM families, human expert benchmarking, and longitudinal tracking of AI generation capabilities as models improve. The SAGE framework demonstrates that systematic theory-driven prompting can direct LLM capabilities toward reliable interpretive evaluation while also revealing the boundaries of current AI literary generation.
